%% fig-sequence.tex -- Figure: end-to-end sequence for one prompt refinement,
%% from a flagged production trace to a rotated alias and its next trace.
%%
%% Seven lifelines. Activation bars show which component holds the work. The two
%% human decisions are named in the message labels, because "how many people are
%% in this loop, and where" is the question the architecture actually answers.
%%
%% Geometry note: lifelines start at x=2.55 to reserve a left margin for the
%% phase braces, and the rightmost lifeline sits at 16.55 so the widest header
%% node still lands inside the 7.0in text block.

\begin{figure}[tbp]
\centering
\begin{cbwide}
\begin{tikzpicture}[x=1cm,y=1cm]

  \def\BOT{-9.86}

  % ---- lifelines -----------------------------------------------------------
  \foreach \x/\s/\n in {%
      2.55/extern/{Production\\agent},
      4.88/surface/{Operator\\(browser)},
      7.22/surface/{Route\\handler},
      9.55/async/{Refinement\\worker},
      11.88/control/{Optimizer +\\EvalProvider},
      14.22/govern/{Governance\\substrate},
      16.55/extern/{\mlflow{}}}
  {
    \node[\s, anchor=north, text width=2.00cm, minimum height=0.60cm,
          font=\figB\bfseries] at (\x,0) {\n};
    \draw[lifeline] (\x,-0.80) -- (\x,\BOT);
  }

  % ---- a message: from-x, to-x, y, style, label ----------------------------
  \newcommand{\msgline}[5]{%
    \draw[#4] (#1,#3) -- (#2,#3);
    \node[elab, anchor=south, font=\figB] at ({(#1+#2)/2},{#3+0.04}) {#5};}

  % ---- activation bars: who is holding the work ---------------------------
  \fill[activation] (9.42,-2.81) rectangle (9.68,-5.87);
  \fill[activation] (11.75,-3.64) rectangle (12.01,-4.96);
  \fill[activation] (7.09,-7.06) rectangle (7.35,-8.72);
  \fill[activation] (14.09,-7.43) rectangle (14.35,-8.72);

  % ---- the sequence --------------------------------------------------------
  \msgline{2.55}{16.55}{-1.19}{msg}{1.~the agent emits a trace during normal
    operation}
  \msgline{16.55}{4.88}{-1.70}{retmsg}{2.~a reviewer flags the response, writing
    an \mlflow{} assessment}

  \msgline{4.88}{7.22}{-2.23}{msg}{3.~\mode{Verify} --- \textbf{human decision}}
  \msgline{7.22}{9.55}{-2.76}{msg}{4.~enqueue a refinement job as a durable row}

  \msgline{9.55}{16.55}{-3.27}{msg}{5.~assemble evidence --- traces,
    assessments, captured examples}
  \msgline{9.55}{11.88}{-3.69}{msg}{6.~diagnose, then run the optimizer}
  \msgline{11.88}{16.55}{-4.13}{msg}{7.~score the candidate against the baseline
    corpus}
  \msgline{16.55}{11.88}{-4.54}{retmsg}{8.~per-dimension scores}
  \msgline{11.88}{9.55}{-4.96}{retmsg}{9.~candidate + scores}

  \msgline{9.55}{14.22}{-5.47}{msg}{10.~submit to the \advgate{}}
  \msgline{14.22}{9.55}{-5.93}{retmsg}{11.~\textsc{enforced} verdict --- advance
    to \code{candidate\_ready}, or stop here}

  \msgline{9.55}{4.88}{-6.51}{retmsg}{12.~the candidate surfaces for review ---
    diff, scores, diagnosis}
  \msgline{4.88}{7.22}{-7.06}{msg}{13.~\mode{Apply} --- \textbf{human decision}}

  \msgline{7.22}{14.22}{-7.52}{msg}{14.~authorize (\operator{} or \admin{}), then
    write the audit row in the same transaction}
  \msgline{14.22}{16.55}{-8.01}{msg}{15.~read the outgoing target, rotate,
    \emph{then} commit}

  \msgline{14.22}{7.22}{-8.72}{retmsg}{16.~rollback checkpoint written --- the
    release is now reconstructable from records}

  % ---- the loop closes -----------------------------------------------------
  \draw[feedback] (16.55,-9.04) -- (16.55,-9.60) -- (2.55,-9.60) -- (2.55,-9.24);
  \node[anchor=south, font=\figB, text=cbAsync, inner sep=1.8pt] at (9.55,-9.58)
    {17.~the next production call resolves \prodalias{} and gets the new version
     --- no client deployment, and its trace is the next signal};

  % ---- phase braces in the reserved left margin ---------------------------
  \foreach \ya/\yb/\lab in {%
      -0.94/-1.98/{observe},
      -2.02/-2.72/{verify},
      -2.76/-5.20/{refine},
      -5.24/-6.72/{gate},
      -6.76/-8.92/{release}}
  {
    \draw[decorate, decoration={brace, amplitude=2.6pt}, draw=cbMuted,
          line width=0.4pt] (1.42,\ya) -- (1.42,\yb);
    \node[anchor=east, font=\figB\itshape, text=cbMuted, inner sep=1.5pt]
      at (1.30,{(\ya+\yb)/2}) {\lab};
  }

\end{tikzpicture}
\end{cbwide}
\caption{End-to-end sequence for one prompt refinement, from a flagged production
trace to a rotated alias and the trace that closes the loop. Activation bars mark
which component holds the work; the left-hand braces group messages into the
phases of \figref{fig:chain}. Three properties are meant to be read off this
figure directly. First, exactly two steps are human decisions --- step~3 confirms
the failure is real and step~13 is an explicit operator action over the diff and
the scores. Nothing between them requires a person, and nothing after step~13
asks for one. Second, the enforced gate at step~11 sits \emph{before} the
candidate reaches review, so an operator is never asked to adjudicate a candidate
that failed its own evaluation. Third, step~15 reads the outgoing alias target and
commits a release intent before rotation, so a crash leaves a durable reconciliation
obligation with the exact value rollback must restore (\algref{alg:promote}).}
\label{fig:sequence}
\figkey
\end{figure}
